\documentclass[11pt]{article}
\usepackage{fontspec,amsmath,amssymb,amsthm,hyperref}
\usepackage[margin=1in]{geometry}
\setlength{\emergencystretch}{3em}
\title{Bounded random Hessians act on L\ensuremath{{}^{2}}}
\author{LangevinLean mathematical exposition}
\date{}
\begin{document}\maketitle
\noindent Source module: \texttt{KineticLangevin/RandomOperatorLp.lean}. Compile this proof document with XeLaTeX or LuaLaTeX. The accompanying HTML pairs the mathematical argument with the exact Lean statements and records the scope of the formalization.


A Hessian evaluated at a random initial position is itself random. To use it inside Bochner time integrals, we construct its genuine continuous linear action on square-integrable random vectors. The operator may depend on every coupled variable; independence is unnecessary.


\section{Construct the action by Hölder multiplication}


A measurable operator field A with a deterministic bound ‖A(ω)‖\ensuremath{\le}M belongs to L∞. Evaluation is a continuous bilinear map from operators and vectors to vectors. Mathlib’s Hölder construction therefore sends the actual L∞ operator and an actual L\ensuremath{{}^{2}} vector to an L\ensuremath{{}^{2}} vector, with 1/∞+1/2=1/2. Fixing the operator gives the continuous linear map randomOperatorLp.



\noindent Lean declarations: \texttt{randomOperatorLp}.

\section{Identify its representatives}


The Holder construction supplies an almost-everywhere representative identity. The representative of the L∞ equivalence class is A(ω) almost everywhere, so the constructed L\ensuremath{{}^{2}} output is exactly A(ω)X(ω). This identity preserves the original probability-space coupling and permits pointwise Taylor and truncation inequalities to be applied to the actual output.



\noindent Lean declarations: \texttt{randomOperatorLp\_coe}.

\section{Retain the deterministic operator bound}


Pointwise, ‖A(ω)X(ω)‖\ensuremath{\le}‖A(ω)‖‖X(ω)‖\ensuremath{\le}M‖X(ω)‖. Integrating this norm inequality gives ‖AX‖\ensuremath{{}_{2}}\ensuremath{\le}M‖X‖\ensuremath{{}_{2}}. Taking the supremum over vectors proves the operator norm of the constructed map is at most M. These estimates justify commuting the random Hessian through time integrals and bounding the integrated force correction.



\noindent Lean declarations: \texttt{randomOperatorLp\_apply\_norm\_le}, \texttt{randomOperatorLp\_norm\_le}.

\end{document}
